\fichatitulo{38 · AVALIAÇÃO}{Preprint v1 · 2022}{HELM}
{\small Liang et al. \textit{Holistic Evaluation of Language Models}. Preprint v1 · 2022. \href{https://arxiv.org/pdf/2211.09110v1}{Fonte consultada}.\par}

\campo{Problema e objetivo}
Como avaliar modelos de linguagem além de uma única medida?

\campo{Principal contribuição · síntese interpretativa}
Propõe avaliação ampla e explícita de cenários, métricas e lacunas de cobertura.

\campo{Método / natureza da evidência}
Compara 30 modelos sob condições padronizadas; combina acurácia, calibração, robustez, justiça, viés, toxicidade e eficiência.

\campo{Resultado ou argumento principal}
Expõe compromissos entre modelos e tarefas, dificultando reduzir resultados a um ranking único.

\campo{Excertos-chave · seleção literal}
\begin{itemize}
\excerto{1}{noting what’s missing or underrepresented}{Resumo.}
\excerto{2}{standardized conditions}{Resumo.}
\end{itemize}

\campo{Limitações e análise crítica}
Autores: validade e confiabilidade dos conjuntos não possuem garantia uniforme; limitações de amostragem e sementes afetam significância. Análise da ficha: escolher métricas exige justificar o que representam no uso parlamentar.

\campo{Aproveitamento na dissertação · proposta}
Metodologia: separar dimensões, reportar variabilidade e evitar agregar pontuações incompatíveis sem justificativa.

\registro{Fonte consultada nesta preparação: Resumo e seção 11. Chave da bibliografia: 2023; versão consultada: novembro de 2022. Chave: \texttt{liangEtAl2023helm}.}
